% ============================================================================
% sec_mechanism_evidence.tex — V8: Stronger mechanism evidence
% Addresses external referee Concern 5: mechanism evidence suggestive
% ============================================================================

\subsection{Bid Rotation and Bid Inflation} % V8: Concern 5
\label{sec:bid_rotation}

We present three additional mechanism tests that go beyond the
Bajari--Ye residual analysis.

\paragraph{Winner persistence.}
For each FL firm, we compute the Herfindahl index of winners across
all tenders in which the firm participates. Under cover bidding, FL
firms should operate in markets with \emph{diversified} winners
(the cartel rotates designated winners across tenders), while
non-FL always-losers---who lose by competitive forces---should
face \emph{concentrated} winners (natural market leaders). The data
confirm this prediction: FL firms' winner HHI is 0.178 (14.3
unique winners), significantly lower than non-FL always-losers'
0.303 (5.0 unique winners; $p < 0.001$). FL firms co-participate
with a wider range of winners, consistent with operating in
competitive markets where cartels rotate designated winners rather
than relying on a single dominant firm.

\paragraph{Repeated FL--winner pairs.}
If FL firms are deployed by specific cartels, they should share
unusually many tenders with their cartel's designated winners. Among
38,941 FL--winner pairs, the mean number of shared tenders is 2.54,
with 4,603 pairs sharing $\geq 5$ tenders and 379 pairs sharing
$\geq 20$ tenders (maximum: 177). By comparison, non-FL
loser--winner pairs share only 1.46 tenders on average, with just
494 pairs reaching $\geq 5$. The substantially heavier right tail of
the FL--winner pair distribution---with hundreds of pairs sharing
10+ tenders---is consistent with persistent cartel relationships in
which the same FL firms repeatedly accompany the same winners.

\paragraph{Bid inflation.}
Under cover bidding, FL bids should be strategically inflated above
the winning price (constraint~\ref{eq:cover_above}: $b_\ell > b^*$).
The median FL bid-to-winner ratio is 1.85 (bids 85\% above the
winner), compared to 1.43 for non-FL losers (43\% above). Controlling
for item and year fixed effects, FL bids are 15.4\% higher relative
to the winning price than non-FL losing bids ($p < 0.001$).
Figure~\ref{fig:bid_inflation} shows the bid-ratio distributions:
the FL distribution is shifted rightward and more dispersed,
consistent with cover bids drawn from a distribution above the
cartel's designated winning bid.

\begin{figure}[htbp]
\centering
\includegraphics[width=0.82\textwidth]{images/fig_bid_inflation.pdf}
\caption{Distribution of log(bid / winning price) for FL and non-FL
losing bids. FL bids are shifted rightward (15.4\% higher, $p < 0.001$)
and more dispersed, consistent with strategically inflated cover bids
drawn from a distribution above the designated winning price.}
\label{fig:bid_inflation}
\end{figure}

\paragraph{Assessment.}
The three tests jointly strengthen the mechanism evidence beyond
``consistent with collusion.'' Winner diversification (Test~1) confirms
Proposition~\ref{prop:market_selection}; persistent FL--winner pairs
(Test~2) indicate stable cartel relationships; and bid inflation
(Test~3) shows that FL bids are not merely unlucky competitive bids
but strategically elevated submissions designed to lose.
